
Las siguientes reglas expresan políticas del dominio del problema. Son
independientes de la tecnología con que se implemente el sistema y
condicionan a uno o más requerimientos: un cambio en cualquiera de ellas
obliga a revisar los requerimientos que la referencian, sin alterar la
arquitectura de la solución.

\subsection*{Cobertura y cálculo de rutas}

\begin{longtable}{@{}KQ@{}}
\toprule
ID & \textbf{Regla} \\
\midrule
\endfirsthead
\toprule
ID & \textbf{Regla} \\
\midrule
\endhead
\bottomrule
\endlastfoot
RN-01 & El área de servicio del sistema es el polígono conformado por las Zonas Territoriales 6, 7 y 8 de la alcaldía Gustavo A. Madero. El origen, el destino y la totalidad del trazo de cualquier ruta calculada deben estar contenidos dentro de ese polígono. \\
RN-02 & El sistema ofrece exactamente tres criterios de enrutamiento: \textit{segura}, que minimiza el riesgo acumulado del trayecto; \textit{rápida}, que minimiza la distancia peatonal; y \textit{balanceada}, que optimiza el compromiso entre ambos sujeto al límite de desvío de RN-03. \\
RN-03 & La longitud de la ruta balanceada no debe exceder en más de 30\% la longitud de la ruta rápida para el mismo par origen--destino. Si ninguna alternativa satisface esa condición, se entrega la ruta más cercana posible. \\
RN-04 & El riesgo de cada tramo vial se expresa como un valor normalizado en el intervalo $[0,1]$. \\
RN-05 & El riesgo de cada tramo se pondera según la franja horaria de la consulta. Las franjas consideradas son diurna y nocturna. \\
\end{longtable}

\subsection*{Datos de riesgo y reportes ciudadanos}

\begin{longtable}{@{}KQ@{}}
\toprule
ID & \textbf{Regla} \\
\midrule
\endfirsthead
\toprule
ID & \textbf{Regla} \\
\midrule
\endhead
\bottomrule
\endlastfoot
RN-06 & El catálogo de categorías de condición de riesgo es único para toda la aplicación y está conformado por las variables resultantes del análisis de las variables propuestas en el marco metodológico. \\
RN-07 & Todo reporte ciudadano se registra inicialmente en estado \textit{pendiente}. La cantidad de confirmaciones necesarias para validarlo dependerá de la calificación del autor: si el autor tiene 5 estrellas, el reporte pasa automáticamente a estado \textit{validado}; si tiene 4 estrellas, requiere 2 confirmaciones de usuarios distintos al autor; y si tiene 3 estrellas, requiere 3 confirmaciones de usuarios distintos al autor. Los reportes que no cumplan con las condiciones de validación permanecerán en estado \textit{pendiente}. Únicamente los reportes validados modificarán los pesos del grafo de riesgo. \\
RN-08 & Un reporte validado pierde vigencia a los [Por definir] días de su última confirmación y deja de afectar los pesos del grafo, de modo que el modelo refleje condiciones actuales del entorno. \\
RN-09 & La contribución de los reportes ciudadanos al riesgo de un tramo no puede modificar en más de [Por definir] el valor obtenido del preprocesamiento de variables ambientales e incidencia oficial. \\
RN-10 & El grafo base de riesgo se recalcula a partir del proceso de preprocesamiento con una periodicidad de [Por definir]. Cada recálculo produce una versión identificada del grafo, referida por RNF-01. \\
RN-11 & El catálogo de puntos de ayuda oficiales es administrado por el responsable del sistema. Los usuarios de la aplicación no pueden crear, modificar ni eliminar registros de este catálogo. \\
RN-12 & Solo los usuarios con cuenta verificada pueden registrar, confirmar o desmentir reportes. Un usuario no puede registrar más de 10 reportes en un periodo de 24 horas. \\
\end{longtable}

\subsection*{Cuenta, enlaces y datos personales}

\begin{longtable}{@{}KQ@{}}
\toprule
ID & \textbf{Regla} \\
\midrule
\endfirsthead
\toprule
ID & \textbf{Regla} \\
\midrule
\endhead
\bottomrule
\endlastfoot
RN-13 & Un enlace de seguimiento es válido únicamente durante la navegación que lo originó. Se invalida de forma automática al finalizar o cancelar esa navegación, y de forma manual cuando el usuario lo revoca. Un enlace invalidado no puede reutilizarse para consultar trayectos posteriores. \\
RN-14 & Los enlaces de verificación de correo y de restablecimiento de contraseña son de un solo uso y expiran a los 10 minutos de su emisión. \\
RN-15 & Al eliminar una cuenta, todos los datos personales asociados a ella deberán eliminarse en un plazo máximo de 7 días. Los reportes realizados por la cuenta podrán conservarse únicamente de forma anonimizada, eliminando cualquier información que permita identificar al usuario, con el fin de utilizarlos para la retroalimentación y mejora del sistema. \\
RN-16 & La ubicación del usuario se utiliza exclusivamente para el cálculo de rutas, la navegación, el registro de reportes y el enlace de seguimiento. No se cede a terceros sin consentimiento expreso del usuario. El tratamiento se rige por el aviso de privacidad, conforme a la Ley Federal de Protección de Datos Personales en Posesión de los Particulares. \\
RN-17 & La verificación del correo electrónico es condición previa para el uso de cualquier función de la aplicación distinta del registro, el inicio de sesión y la recuperación de contraseña. \\
RN-18 & Parámetros operativos del sistema: radio de proximidad para notificaciones de nuevos reportes, 1.5 metros; umbral de desvío que dispara el recálculo de la ruta, 50 metros. Ambos son configurables sin necesidad de publicar una nueva versión de la aplicación. \\
\end{longtable}